\section{The full polarized adjoint: a checked application of Thorne's theorem}
\label{sec:adk-match}

This section resolves the matching question called
\texttt{AD\_VANISH\_K\_MATCH} in the v0.41 register.  The conclusion is
the entire original \texttt{AdVanishK}, for every rank $n\ge2$; the
selected Kato application uses rank three.  The proof is a published
theorem application with explicit coefficient, normalization, action,
and local-condition comparisons.  It is independent of the unfinished
ordinary Hecke and host-construction arguments.  The separate
\texttt{HOST\_RIGID\_S} obligation is not closed by this result.

\subsection{The exact statement to be supplied}

Let $K$ be an imaginary quadratic field, let $\Omega$ be an algebraic
closure of $K$ containing the specified copy of $\mathbb Q$, let $\ell$
be any rational prime, and fix
$\iota:\overline{\mathbb Q}_{\ell}\simeq\mathbb C$.
Write $G_K=\operatorname{Gal}(\Omega/K)$,
$G_{\mathbb Q}=\operatorname{Gal}(\Omega/\mathbb Q)$, and
$\epsilon:G_K\to\overline{\mathbb Q}_{\ell}^{\times}$ for the actual
cyclotomic character, obtained by the specified inclusion
$\mathbb Z_\ell\hookrightarrow\overline{\mathbb Q}_\ell$.
Suppose $n\ge2$ and
\[
 r:G_K\longrightarrow\operatorname{GL}_n(\overline{\mathbb Q}_\ell)
 \tag{ADK.1}
\]
is continuous.  The original automorphy hypothesis says that there is
a regular algebraic cuspidal representation $\pi$ of
$\operatorname{GL}_n(\mathbb A_K)$, conjugate self-dual, such that $r$
is semisimple and has the following unramified characteristic
polynomials outside a finite set.  If $w$ has norm $q_w$, if $F_w$ is
\emph{arithmetic} Frobenius, and if
$\{\alpha_{w,1},\ldots,\alpha_{w,n}\}$ is its Satake multiset in the
source normalization, then
\[
 \det\bigl(T-\iota r(F_w^{-1})\bigr)
 =\prod_{i=1}^{n}
   \bigl(T-q_w^{(n-1)/2}\alpha_{w,i}\bigr).
 \tag{ADK.2}
\]
Thus (ADK.2) uses geometric Frobenius; it is not the formula obtained
by applying the same right side to $F_w$.
The remaining irreducibility hypothesis is
\[
                       r|_{\ker\epsilon}\text{ is irreducible}.
 \tag{ADK.3}
\]
It is a characteristic-zero assertion over the full infinite
cyclotomic kernel, not residual irreducibility or irreducibility after
adjoining just $\ell$th roots of unity.

Let $V=M_n(\overline{\mathbb Q}_\ell)$ and let $\rho$ be a
$\overline{\mathbb Q}_\ell$-linear representation of $G_{\mathbb Q}$
on $V$.  The original \texttt{IsAdOf} hypothesis supplies a complex
conjugation $c\in G_{\mathbb Q}$, an embedding
$j:\Omega\hookrightarrow\mathbb C$ inducing that conjugation, and
$P\in\operatorname{GL}_n(\overline{\mathbb Q}_\ell)$, such that
\begin{align}
 r(s)^tP r(csc^{-1})&=\epsilon(s)^{1-n}P &&(s\in G_K),
                 \tag{ADK.4}\\
 \rho(s)X&=r(s)Xr(s)^{-1} &&(s\in G_K),
                 \tag{ADK.5}\\
 \rho(c)X&=-P^{-1}X^tP &&(X\in V).
                 \tag{ADK.6}
\end{align}
The intended conclusion is that every continuous $\rho$-cocycle
$a:G_{\mathbb Q}\to V$ satisfying the following two conditions is a
coboundary:
\begin{enumerate}
\item At every finite $p\ne\ell$, its restriction to inertia is
  $s\mapsto\rho(s)v_p-v_p$ for some $v_p\in V$.
\item At $\ell$, after identifying the local decomposition group
  using any permitted embedding into the fixed local algebraic
  closure, $a(s)\otimes1=\rho_{\mathrm{cris}}(s)b-b$ for some
  $b\in V\otimes_{\mathbb Q_\ell}B_{\mathrm{cris}}$.
\end{enumerate}
Here the action on the tensor product is the diagonal action, including
the Galois action on $B_{\mathrm{cris}}$.  These are the literal
\texttt{IsBK.away} and \texttt{IsBK.above} conditions, as verified
below.  In particular the conclusion is about all of $M_n$, including
its scalar matrices.

% BEGIN ADK1
\begin{lemma}[ADK1: one finite coefficient field for a representation and a cocycle]
\label{claim:ADK1}
A continuous finite-dimensional representation of a profinite group
over $\overline{\mathbb Q}_\ell$ has all its matrix coefficients in
one finite extension of $\mathbb Q_\ell$.  Once its coefficients are
in $E$, each continuous cocycle has all its values in a finite
extension $E'/E$.  Any finitely many fixed matrices may be included
in $E'$.
\end{lemma}
\begin{proof}
NR19 gives a countable cofinal family of finite subfields of
$\overline{\mathbb Q}_\ell$.  Briefly, approximate any algebraic
$\ell$-adic number by an element of a fixed countable dense copy of
$\overline{\mathbb Q}$ more closely than its distance to its other
conjugates.  The nonarchimedean triangle inequality makes every
automorphism fixing the approximant fix the original number.  The
fixed-field theorem then puts that number in the finite local field
generated by the approximant.  Taking finite subsets gives the claimed
cofinal family; taking successive composita makes it increasing.
Each finite local field is complete, hence closed in the algebraic
closure equipped with its $\ell$-adic topology.

For a continuous representation $u$, the inverse images
$u^{-1}(\operatorname{GL}_d(E_i))$ are closed subgroups covering the
compact group.  Baire's theorem makes one of them have nonempty
interior, hence be open.  It has finite index.  Adjoining the entries
of $u$ at finitely many coset representatives captures every matrix
coefficient in a finite field.  Adjoin the entries of any additional
fixed matrices as well.

If the action is already defined over $E$, enlarge every $E_i$ to
contain $E$.  For a cocycle $a$, the sets
$H_i=\{g:a(g)\in E_i^d\}$ are again closed subgroups: the identities
\[
 a(gh)=a(g)+u(g)a(h),\qquad
 a(g^{-1})=-u(g^{-1})a(g)
\]
show closure under product and inverse.  Apply the same Baire and
finite-coset argument.  This argument does not assume that an arbitrary
compact subset of the algebraic closure is contained in a finite
extension; it uses the group or cocycle laws essentially.
\end{proof}
% END ADK1

% BEGIN ADK2
\begin{lemma}[ADK2: the exact infinite cyclotomic subgroup]
\label{claim:ADK2}
With the original embeddings,
\[
 \ker\epsilon=G_{K(\mu_{\ell^\infty})}.
 \tag{ADK.7}
\]
Consequently (ADK.3) is exactly the irreducibility hypothesis in
Thorne's Theorem~6.1.
\end{lemma}
\begin{proof}
The $\ell$-adic cyclotomic character is the compatible system of
actions on $\mu_{\ell^m}$, for all $m\ge1$.  The inclusion of
$\mathbb Z_\ell$ into the coefficient algebraic closure is
injective.  Thus $\epsilon(s)=1$ if and only if $s$ fixes every such
root of unity, which is equivalent to fixing their compositum with
$K$.  The induced topology is the closed-subgroup topology on both
sides.  This identifies the actual infinite subgroup, without taking
a finite cyclotomic approximation.  Irreducibility on this subgroup
also implies irreducibility on $G_K$.
\end{proof}
% END ADK2

% BEGIN ADK3
\begin{lemma}[ADK3: the normalized attached representation is the given one]
\label{claim:ADK3}
The automorphic representation in (ADK.2) is polarizable, and its
standard normalized Galois representation $r_{\pi,\iota}$ is
isomorphic to the given $r$.  The polarization has multiplier
$\epsilon^{1-n}$ on $G_K$ and a symmetric pairing at the specified
complex conjugation.
\end{lemma}
\begin{proof}
We first unpack the automorphic dictionary; no conclusion is inferred
from a predicate's name.  In
\texttt{AutomorphicMembership/ExpoConjOfIsCsd.lean},
\texttt{CuspSub} is a nonzero irreducible submodule of smooth
functions on $\operatorname{GL}_n(\mathbb A_K)$, stable under the
finite adelic group, the maximal compact group at infinity, and the
complexified Lie algebra.  Its vectors are left
$\operatorname{GL}_n(K)$-invariant, have an open finite-adelic
stabilizer, are finite under the maximal compact and the center of
the enveloping algebra, and have moderate growth.  The literal
\texttt{IsCuspidal} tests the maximal standard parabolics indexed
by $0<r<n$.  Vanishing for every proper standard parabolic follows
by integrating first over a containing maximal parabolic's
unipotent radical and then over the remaining quotient; the
unipotent rational quotients are compact, so Fubini applies.
Thus the integral of each vector over every proper standard
parabolic unipotent radical is zero.  These are the usual cuspidal
automorphic $(\mathfrak g,K_\infty)\times G(\mathbb A_f)$ modules;
the standard smooth, admissible realization of an irreducible
cuspidal automorphic representation is precisely this realization.
\texttt{CG.mk} passes to its intertwining-isomorphism class and
changes none of these actions.

There are two normalization checks in this dictionary.
First, the source's central polynomial is
\[
 \operatorname{cap}_{\tau}(u)
    =\operatorname{qdet}(E_{\tau},u+(n-1)/2).
\]
On a highest-weight vector of weight $(\lambda_0,\ldots,
\lambda_{n-1})$ it has eigenvalue
$\prod_{i=0}^{n-1}(u+\lambda_i+(n-1)/2-i)$;
this is the literal \texttt{cap\_highest} formula in
\texttt{AutomorphicMembership/ContinuousConjFin.lean},
lines 147--157.  Thus its roots record the ordinary
Harish--Chandra parameters with the usual half-sum-of-positive-roots
shift.  \texttt{CG.IsRegAlg}, through
\texttt{isRegAlg\_iff}, requires for every embedding $\tau$ that
these $n$ parameters are distinct and belong to
$\mathbb Z+(n-1)/2$.  Sorting them as
$z_0>\cdots>z_{n-1}$ gives the integral dominant weight
$\lambda_i=z_i-(n-1)/2+i$: indeed
$\lambda_i-\lambda_{i+1}=z_i-z_{i+1}-1\ge0$.
Conversely every integral dominant weight produces such distinct
parameters.  This is exactly regular algebraicity in BLGGT,
\S2.1, for both complex embeddings.  The source defines a fallback
multiset consisting of $n$ zeros when an infinitesimal character
does not exist; since $n\ge2$, distinctness excludes that fallback.
Consequently the regularity hypothesis really does supply the
stated central eigencharacter.

Second, \texttt{CG.HasSat} in
\texttt{GaloisCohomology2/StrongMultiplicityOne.lean},
lines 61--66, descends the predicate
\texttt{HasSatGL} of the actual local representation.  With
$\operatorname{vol}(\operatorname{GL}_n(\mathcal O_w))=1$, it
says that the spherical fixed space is nonzero and that, for every
nonzero fixed vector and every $0\le j\le n$, the characteristic
function of the double coset of
$\operatorname{diag}(\varpi I_j,I_{n-j})$ acts by
\[
             q_w^{j(n-j)/2} e_j(\alpha_{w,1},\ldots,
                                    \alpha_{w,n}).
 \tag{ADK.9}
\]
This is obtained by unfolding \texttt{IsSatVecGL} and
\texttt{satEig}, not by selecting a different local constituent.
The factor in (ADK.9) is exactly that of normalized Satake
parameters, or equivalently normalized unramified induction.
The elementary symmetric functions determine the multiset, so the
source parameters and the parameters used by normalized local
Langlands agree.  The source's further \texttt{algSat} multiplies
each of them by $q_w^{(n-1)/2}$, as displayed in (ADK.2).

Finally \texttt{DualIsConj} is a nondegenerate complex bilinear
pairing $B$ on the cusp module with
$B(\pi(g)v,\pi(g^c)w)=B(v,w)$ for finite adelic and compact
elements, and
$B(Xv,w)+B(v,X^cw)=0$ for Lie algebra elements.
Here \texttt{conjOp} is entrywise CM conjugation at finite places,
and the corresponding conjugate actions on compact and Lie algebra
elements.  The map $v\mapsto B(v,\cdot)$ intertwines $\pi$ with
the contragredient of $\pi^c$.  It lands in the smooth,
$K_\infty$-finite contragredient: a compact open subgroup fixing
$v$ fixes its functional by the first pairing identity, and the
$K_\infty$-span of that functional is the image of the finite
$K_\infty$-span of $v$.  Nondegeneracy and irreducibility make this
an isomorphism of admissible irreducible modules.  This checks
$\pi^c\simeq\pi^\vee$, with a bilinear pairing rather than an
unintended Hermitian one.

Let $\delta_K$ denote the quadratic Hecke character of
$K/\mathbb Q$.  Its pullback by the idele norm from $K$ is trivial,
and its infinity value at $-1$ is $-1$.  Since
$\pi^c\simeq\pi^\vee$, the pair
$(\pi,\delta_K^n)$ is polarized with the precise infinity sign
$(-1)^n$ required in BLGGT, \S2.1.  Regular algebraicity and
cuspidality are the original hypotheses, not new conclusions of a
finite coefficient argument.

BLGGT, Theorem~2.1.1, attaches a continuous semisimple
$r_{\pi,\iota}$ and identifies its local Weil--Deligne representation
away from $\ell$ with that of
$\pi_w\otimes|\det|_w^{(1-n)/2}$.  At an unramified place the
geometric Frobenius eigenvalues are therefore exactly
$q_w^{(n-1)/2}\alpha_{w,i}$.  This is (ADK.2), with the same
$\iota$ and the same direction of Frobenius.

By ADK1 put both continuous representations over one finite
coefficient field.  Chebotarev density makes equality of their
unramified Frobenius traces outside a finite set imply equality of
their continuous traces on $G_K$.  The characteristic-zero
semisimple character criterion then supplies an isomorphism.  One
elementary proof of that criterion uses the algebra generated by the
two actions on their direct sum: its radical kills a semisimple
module, and its simple matrix factors have independent trace
characters.  Equality of traces gives equal multiplicities of all
simple factors.  Thus we may transport $r_{\pi,\iota}$ to the
original space carrying $r$.

The Galois multiplier in BLGGT is
$\delta_K^n\epsilon^{1-n}$ on $G_{\mathbb Q}$.  Its restriction to
$G_K$ is $\epsilon^{1-n}$; its value at complex conjugation is
$(-1)^n(-1)^{1-n}=-1$.  The total oddness assertion in that theorem
and the polarization convention in \S2.1 therefore give a symmetric
nondegenerate matrix $P_0$ with
$r(s)^tP_0r(csc^{-1})=\epsilon(s)^{1-n}P_0$.
Every complex conjugation arising from an embedding of $\Omega$ in
$\mathbb C$ is allowed in this assertion, including the one in
(ADK.4).  The quadratic character used for the infinity sign is
trivial on $G_K$ and changes neither the attached representation nor
its Frobenius polynomials.
\end{proof}
% END ADK3

% BEGIN ADK4
\begin{lemma}[ADK4: equality of the full polarized adjoint actions]
\label{claim:ADK4}
The original $\rho$ in (ADK.5)--(ADK.6) is the full unitary adjoint
$\operatorname{ad}r_{\pi,\iota}$ in Thorne's theorem.
\end{lemma}
\begin{proof}
The matrices $P$ and $P_0$ satisfy the same intertwining equation.
Dividing their equations shows that $P^{-1}P_0$ commutes with
$r(cG_Kc^{-1})=r(G_K)$.  This representation is irreducible over
the algebraically closed coefficient field by ADK2, so Schur's lemma
makes $P_0$ a nonzero scalar multiple of $P$.  Hence
\[
              -P_0^{-1}X^tP_0=-P^{-1}X^tP.
 \tag{ADK.8}
\]
For clarity, the minus-transpose formula is the differential of the
disconnected component of the standard group
$\mathcal G_n=(\operatorname{GL}_n\times\operatorname{GL}_1)
\rtimes\{1,j\}$, where
$j(g,a)j^{-1}=(a g^{-t},a)$; see BLGGT, \S1.1.
The differential of $g\mapsto g^{-t}$ at the identity is
$X\mapsto-X^t$.  With the pairing $P_0$, conjugation by the
corresponding component gives $X\mapsto-P_0^{-1}X^tP_0$.
The connected subgroup acts by $X\mapsto r(s)Xr(s)^{-1}$.
Thus the two actions agree on $G_K$ and on $c$.  Since $K$ is
imaginary quadratic, those elements generate $G_{\mathbb Q}$,
proving equality on the entire original group.

Thorne's adjoint acts on $\operatorname{Lie}\operatorname{GL}_n=M_n$;
it is not restricted to trace-zero matrices.  In particular $c$
acts by $-1$ on the scalar line.  No trace-zero substitution, omitted
quadratic line, or division by $n$ is used.  The action is continuous:
on each of the two open cosets of $G_K$ it is expressed by the
continuous matrices of $r$ and the fixed matrix $P$.
\end{proof}
% END ADK4

% BEGIN ADK5
\begin{lemma}[ADK5: the original Bloch--Kato condition, including coefficients]
\label{claim:ADK5}
For the representation $\rho$ of ADK4, the source's
\texttt{BlochKatoSelmerVanishes} is precisely vanishing of the usual
Bloch--Kato Selmer group over $\mathbb Q$, with its literal
algebraic-closure coefficient and cocycle scope.
\end{lemma}
\begin{proof}
The original \texttt{Z1} consists of continuous cocycles.  At
$p\ne\ell$, the assertion that a cocycle is a coboundary on inertia
is exactly membership in the unramified local subgroup
$\ker(H^1(G_{\mathbb Q_p},V)\to H^1(I_p,V))$.
Changing the stored valuation above $p$ conjugates the decomposition
and inertia groups.  For a global cocycle the change of restriction
is conjugation followed by a coboundary, so this kernel condition is
unchanged.  Thus the existential choice of a place in
\texttt{IsBK.away} is the usual local condition.

At $\ell$, \texttt{decJ} is the graph of the restriction map from
the local absolute Galois group along the chosen embedding
$\Omega\hookrightarrow\overline{\mathbb Q}_\ell$.  The second
projection identifies this graph with the local group.  Indeed the
image of $\Omega$ is the subfield of elements algebraic over the
specified global field, which is stable under local automorphisms;
restriction determines the first component uniquely.  This
identification is topological, too.  A basic open subgroup of the
global Galois group fixes a finite set $S\subset\Omega$.
Its inverse image under local restriction fixes the finite extension
$\mathbb Q_\ell(j(S))$, and is therefore open in the local Galois
group.  Restriction is continuous; hence the inverse to the second
projection, $s\mapsto(\operatorname{res}_j(s),s)$, is continuous.
The second projection is itself continuous as a restricted product
projection.  In particular a continuous global cocycle pulls back
to the usual continuous local cocycle.  On the graph,
\texttt{actC} is precisely the diagonal action on
$V\otimes_{\mathbb Q_\ell}B_{\mathrm{cris}}$.
The original \texttt{finiteCocycles} requires a tensor primitive,
which says exactly that the local cohomology class dies after this
tensor extension.  BX1 identifies the original tied period-ring
model with the usual $B_{\mathrm{cris}}$, equivariantly and with
the stated cyclotomic character.  This proves the local comparison
without replacing finite cohomology by a larger de Rham condition.
At the real place positive-degree cohomology is zero by averaging
over a group of order two in characteristic zero.

Finally, ADK1 gives a finite field $E$ containing all coefficients
of the action and all values of an individual continuous cocycle.
Every local algebraic-closure primitive uses finitely many vector
coefficients, so lies over some finite $E'/E$.  An $E$-linear
retraction $q:E'\to E$ with $q(1)=1$ commutes with the $E$-defined
action.  Apply it to an inertia primitive, or apply
$q\otimes1$ to a crystalline tensor primitive, to return it to
$E$.  Explicitly, for the coordinatewise retraction $q_V$, a pure
tensor $v\otimes b$ and a local element $s$ satisfy
\[
 (q_V\otimes1)\bigl(\rho(s)v\otimes s(b)\bigr)
     =\rho(s)q_V(v)\otimes s(b)
     =s\bigl((q_V\otimes1)(v\otimes b)\bigr).
\]
Linearity extends this equality to every tensor; it uses the
$E$-linearity of $q_V$ and the identity on the period-ring factor,
so includes the nontrivial $B_{\mathrm{cris}}$ action.  The same map
fixes $a(s)\otimes1$ since $a(s)\in V_E$.
This is done separately at each place; it does not posit one
finite field containing infinitely many selected local primitives.
Therefore the source cocycle belongs to $H^1_f(\mathbb Q,V_E)$.
Conversely a finite-coefficient Selmer cocycle remains one after
scalar extension.  A coboundary in either presentation has the
literal form $g\mapsto\rho(g)v-v$, as required by the source.
\end{proof}
% END ADK5

% BEGIN ADK6
\begin{theorem}[ADK6: the full original \texttt{AdVanishK}]
\label{claim:ADK6}
Under (ADK.1)--(ADK.6), every source Bloch--Kato cocycle is a
coboundary.  This proves \texttt{AdVanishK} with all its original
quantifiers, in particular for rank three and for every prime
allowed there.
\end{theorem}
\begin{proof}
We use Thorne, \emph{On the vanishing of adjoint Bloch--Kato Selmer
groups of irreducible automorphic Galois representations},
Theorem~6.1 (arXiv:2207.04925v1, printed p.~24): a regular algebraic
cuspidal polarized representation over a CM field has vanishing
unitary adjoint Selmer group when its attached representation
remains irreducible over the infinite cyclotomic extension.
Apply it to $F=K$, $F^+=\mathbb Q$, the same $n$, prime $\ell$ and
embedding $\iota$, and the polarized pair of ADK3.  Its
irreducibility hypothesis is exactly ADK2.  Its attached
representation and full adjoint are the original ones by ADK3--ADK4.
Thus $H^1_f(\mathbb Q,\rho)=0$.  ADK5 transfers this statement to
every literal source cocycle, proving the claimed endpoint.

There is no residual-image assumption, bound $\ell>n$, exclusion of
$\ell=3$, ordinary hypothesis, or hypothesis that every place splits
in $K$ in the cited theorem.  The split/unramified/Iwahori conditions
introduced later inside Thorne's proof are reductions in that proof,
not extra assumptions of Theorem~6.1.  This is why the application
does not import the much larger ordinary-lifting package used by a
different source route.
\end{proof}
% END ADK6

\subsection{The separate host rigidity question}

The remaining host endpoint is now a located reconstruction task,
not an unidentified reference.  The exact source is
\texttt{EHostRigidU.hostRigidUE\_alt}; its final assembly occupies
lines 390--440.  Its conclusion concerns the original quotient
$\Gamma=G_{\mathbb Q,S}$, the original marked element $g_1$, a point
which becomes integral after the specified torus conjugation, and
every dual-number deformation satisfying its prescribed local flags.
The source representation is the full adjoint of
$r:G_{\mathbb Q,S}\to\operatorname{GL}_3\rtimes\{\pm1\}$, pulled
back by the actual quotient map.

Its \texttt{BPrime\_s.cyc} already states irreducibility on
$\ker(\eta_4)\cap\ker(\epsilon)$, and
$\ker(\eta_4)=G_{\mathbb Q(i)}$.  Its automorphy field supplies the
cyclotomic twist $r|_{G_{\mathbb Q(i)}}\otimes\epsilon^{-1}$.
The twist has trivial conjugate-pair product after its prescribed
multiplier adjustment, and scalar twists do not change the connected
adjoint action.  A full matching proof must still write the
disconnected action on the same pairing.  More importantly, global
Selmer vanishing yields the host's rigidity only after showing that
each allowed tangent cocycle satisfies the original local finite
conditions.

The away-from-$\ell$ part is finite-index averaging and the literal
diagonal-conjugacy condition.  The substantive local step is
\texttt{rig\_above}, lines 145--191 of \texttt{EHostRigidU.lean}:
it uses \texttt{rig\_sub} and \texttt{rig\_eig} to obtain an actual
open subgroup, fixed inertia weights $(e,0,-e)$, and the required
non-cyclotomic eigenvalue ratios before applying
\texttt{RigLoc.rig\_loc}.  Those hypotheses are not automatic from
the phrase ``ordinary''.  \texttt{RigFin.rig\_fin} and
\texttt{RigDesc.rig\_desc} give the finite-coefficient local
crystalline-extension and descent arguments.  Once they are written
for these same objects, the quotient map's surjectivity descends the
global coboundary back to the original tangent cocycle; its sign must
agree with the source's \texttt{cob} convention.

Accordingly ADK6 discharges \texttt{AD\_VANISH\_K\_MATCH}, subject to
integration review.  It does not discharge \texttt{HOST\_RIGID\_S},
\texttt{RIG\_II}, ordinary support, or the existence of any
automorphic host.  The suggested next host packet is the complete
\texttt{rig\_sub}--\texttt{rig\_eig}--\texttt{rig\_loc}
comparison plus the explicit twist/pairing identification.

\subsection{References and literal source correspondence}

The source pin throughout is
\texttt{7cbf7a0c32fca1badf05931879a37c3f8a2b03b4}.
\texttt{AutomorphicForms/CongruenceNumberRankThree.lean},
lines 158--175, gives \texttt{IsAdOf} and \texttt{AdVanishK}.
\texttt{AutomorphicForms/BlochKatoSelmerVanishes.lean},
lines 69--86 and 117--144, gives the normalized attachment,
automorphy and exact local Selmer definitions.
\texttt{PadicHodgeTheory/FiniteCohomology.lean},
lines 122--151, gives the tensor action and finite cocycles.
These paths are relative to
\texttt{src/Development/AutomorphicArithmetic/AutomorphicVarieties/}.
The accompanying JSON records the full paths, ranges and hashes.

The published statements used are:
\begin{enumerate}
\item J.~A. Thorne, \emph{On the vanishing of adjoint Bloch--Kato
  Selmer groups of irreducible automorphic Galois representations},
  Pure Appl. Math. Q. \textbf{18} (2022), no.~5, Theorem~6.1;
  \url{https://arxiv.org/html/2207.04925v1} and
  \url{https://arxiv.org/pdf/2207.04925}.  The full adjoint is
  specified on printed pp.~3--4; the theorem is on p.~24.
\item T.~Barnet-Lamb, T.~Gee, D.~Geraghty and R.~Taylor,
  \emph{Potential automorphy and change of weight},
  Ann. of Math. \textbf{179} (2014), 501--609.
  Section~1.1, pp.~512--513, defines $\mathcal G_n$;
  Section~2.1, p.~535, specifies the pairing and sign;
  Theorem~2.1.1, p.~538, supplies normalized attachment and
  total oddness.
  \url{https://annals.math.princeton.edu/wp-content/uploads/annals-v179-n2-p03-p.pdf}.
\item Chebotarev density and the characteristic-zero semisimple
  character criterion, used with both representations in one finite
  coefficient field.  The Chebotarev application is the same standard
  one retained in RX5; the character argument is recalled in ADK3.
\end{enumerate}
